% ---- pg-399 （书页 387） ----

残差矩阵为
\[
\mathbf{R}-\hat{\mathbf{L}}\hat{\mathbf{L}}'-\hat{\boldsymbol{\Psi}}=\begin{bmatrix}
0&0.001&-0.002&0.000&0.052\\
0.001&0&0.002&0.000&-0.033\\
-0.002&0.002&0&0.000&0.001\\
0.000&0.000&0.000&0&0.000\\
0.052&-0.033&0.001&0.000&0
\end{bmatrix}
\]
$\mathbf{R}-\hat{\mathbf{L}}\hat{\mathbf{L}}'-\hat{\boldsymbol{\Psi}}$ 的元较对应于主成分因子化 $\mathbf{R}$ 的残差阵的元要小得多, 后者列在例 9.4 中. 有鉴于此, 我们更欣赏极大似然法, 并且在以后的例子中有代表性地描绘它的特点.

由因子解释的占样本总方差的累积比, 主成分因子化要比极大似然因子化大. 这个判别准则典型地倾向于主成分因子化, 这一点并不奇怪. 由主成分因子分析得到的载荷与主成分有关, 它有目的地使方差得到优化(见\hyperref[eq:8-19]{式(8-19)}前的讨论).

把注意力集中到极大似然解, 我们看到所有变量在 $F_1$ 上有大的正载荷. 我们称这个因子为市场因子, 一如我们在主成分分解中所做的那样. 然而, 第二因子的阐述就不像它在主成分分解中的那样清楚了. 银行股有大的正载荷而石油股的载荷可以忽略. 由此看来, 第二因子区分银行股与石油股, 由此大概可以解释为行业因子. 另一种说法是: 第二因子为银行因子. 极大似然解的初始因子载荷的模式受到 $\hat{\mathbf{L}}'\hat{\boldsymbol{\Psi}}^{-1}\hat{\mathbf{L}}$ 为对角阵这个唯一性条件的制约. 因此, 常常直到完成因子旋转时, 有用的因子模式才被发现(见 9.4 节). \hfill$\blacksquare$

\noindent\textbf{例 9.6}\quad (奥林匹克十项全能数据的因子分析)

林登(Linden, 见文献[11])最初对第二次世界大战后到 70 年代中期的奥林匹克十项全能的所有 160 个完整的成绩作过因子分析. 按照他的方法, 我们检查从 1960 年到 2004 年奥林匹克十项全能的所有 280 个完整的成绩. 对(十项全能的)每项比赛的记录值作了标准化, 计时的项目的符号作了改变, 使得所有项目的得分越大越好. 基于 280 个记录, 我们分析协方差矩阵
$\mathbf{R}=$
{\small
\[
\begin{bmatrix}
1.000&0.638\ 6&0.475\ 2&0.322\ 7&0.552\ 0&0.326\ 2&0.350\ 9&0.400\ 8&0.182\ 1&-0.035\ 2\\
0.638\ 6&1.000\ 0&0.495\ 3&0.566\ 8&0.470\ 6&0.352\ 0&0.399\ 8&0.516\ 7&0.310\ 2&0.101\ 2\\
0.475\ 2&0.495\ 3&1.000\ 0&0.435\ 7&0.253\ 9&0.281\ 2&0.792\ 6&0.472\ 8&0.468\ 2&-0.012\ 0\\
0.322\ 7&0.566\ 8&0.435\ 7&1.000\ 0&0.344\ 8&0.350\ 3&0.365\ 7&0.604\ 0&0.234\ 4&0.238\ 0\\
0.552\ 0&0.470\ 6&0.253\ 9&0.344\ 9&1.000\ 0&0.154\ 6&0.210\ 0&0.421\ 3&0.211\ 6&0.412\ 5\\
0.326\ 2&0.352\ 0&0.281\ 2&0.350\ 3&0.154\ 6&1.000\ 0&0.255\ 3&0.416\ 3&0.171\ 2&0.000\ 2\\
0.350\ 9&0.399\ 8&0.792\ 6&0.365\ 7&0.210\ 0&0.255\ 3&1.000\ 0&0.403\ 6&0.417\ 9&0.010\ 9\\
0.400\ 8&0.516\ 7&0.472\ 8&0.604\ 0&0.421\ 3&0.416\ 3&0.403\ 6&1.000\ 0&0.315\ 1&0.239\ 5\\
0.182\ 1&0.310\ 2&0.468\ 2&0.234\ 4&0.211\ 6&0.171\ 2&0.417\ 9&0.315\ 1&1.000\ 0&0.098\ 3\\
-0.035\ 2&0.101\ 2&-0.012\ 0&0.238\ 0&0.412\ 5&0.000\ 2&0.010\ 9&0.239\ 5&0.098\ 3&1.000\ 0
\end{bmatrix}
\]
}
从主成分因子分析的角度看, $\mathbf{R}$ 的前 4 个特征值 4.21, 1.39, 1.06, 0.92 提示因子个数解为 $m=3$ 或 $m=4$. 后续的解释与林登的最初分析非常类似, 否定了 $m=4$ 的解.

在此案例中, 两种解法产生非常不同的结果. 对于主成分因子分析, 除去 1 500 米跑外, 其他的项目在第一因子上都有大的正载荷. 这个因子应该识别为一般运动能力. 因子 2 在 400 米跑和 1 500 米跑两个项目上有非常大的载荷, 因此可以识别为跑的持续力因子. 其他因子不易解释.
